Dans une expérience de type Millikan, une gouttelette d'huile de masse \qty{\v{m}e-15}{kg} reste immobile entre deux plaques horizontales lorsque le champ électrique vertical vaut \qty[round-mode=figures,round-precision=4]{\fpeval{\v{m}*1e-15*9.81/(\v{n}*1.602e-19)}}{N/C}.

\begin{enumerate}[label=\alph*)]
    \item Faites le diagramme des forces sur la gouttelette. Le champ doit-il pointer vers le haut ou vers le bas si la gouttelette est chargée négativement ?
    \item Déterminez la charge de la gouttelette.
    \item Combien d'électrons en excès porte la gouttelette ($e = \qty{1.602e-19}{C}$) ?
\end{enumerate}
